%! Modeling with Quadratics — Unit 2, Lesson 2.7 — Homework
\documentclass[10pt]{article}
\usepackage{algebra2-article}
\usepackage{algebra2-boxes}
\newcommand{\TallMath}[1]{$\displaystyle #1\rule[-1.4em]{0pt}{3.2em}$}

\begin{document}

\pageheader{Unit 2, Lesson 2.7}{Homework}

\vspace{0.10in}

\begin{notesbox}{Practice}
Show your work. \textbf{Build the model, pick the feature the question wants, and answer in context with units.}

\begin{enumerate}[leftmargin=1.9em, itemsep=9pt]
  \item \textbf{Projectile.} A ball is thrown up at $32$ ft/s from a $128$-ft ledge: $h(t)=-16t^2+32t+128$.
    \begin{enumerate}[label=(\alph*), leftmargin=1.9em, itemsep=3pt]
      \item Starting height: \blank{2.4cm}
      \item Maximum height and the time it occurs: \blank{3.6cm}
      \item When it hits the ground: \blank{2.6cm}
    \end{enumerate}

  \item \textbf{Maximum area (three sides).} A rectangular pen against a barn wall uses $48$ ft of fence on three
        sides ($x=$ width). Write the area model, then find the width, the maximum area, and the dimensions.
    \begin{work}
      A(x) &= x(48-2x) = -2x^2+48x \\
      x &= -\frac{48}{2(-2)} = 12\text{ ft} \\
      A(12) &= 12(24) = 288\text{ ft}^2 \quad (12\times24)
    \end{work}

  \item \textbf{Maximum revenue.} A gym has $500$ members paying \$20/month. A study says each \$1 increase loses
        $10$ members ($x=$ number of \$1 increases). Write $R(x)$, then find the price and the maximum revenue.
    \begin{work}
      R(x) &= (20+x)(500-10x) = -10x^2+300x+10000 \\
      x &= -\frac{300}{2(-10)} = 15 \\
      \text{price} &= 20+15 = \$35 \\
      R(15) &= \$12{,}250
    \end{work}

  \item \textbf{Interpret the features.} A company's monthly revenue is $R(x)=-20x^2+40x+1600$ (dollars, $x=$ \$1
        price increases). State what each feature means \emph{in context}:
    \begin{enumerate}[label=(\alph*), leftmargin=1.9em, itemsep=3pt]
      \item the $y$-intercept $(0,1600)$: \blank{5.2cm}
      \item the vertex $(1,1620)$: \blank{5.2cm}
    \end{enumerate}

  \item \textbf{Find the error.} Dana models a pen against a wall ($40$ ft of fence, three sides) as $A(x)=x(40-x)$
        and gets a maximum at $x=20$. What did Dana set up wrong, and what is the correct maximum area?
        \writeline
\end{enumerate}
\end{notesbox}

\begin{extensionbox}
\textbf{1. Same height (systems tie-back).} Two drones: $h_A(t)=-16t^2+64t$ and $h_B(t)=32t$. Other than at launch,
when are they at the same height, and what height is that?
\[ t=\blank{1.6cm} \qquad \text{height}=\blank{2.0cm} \]

\textbf{2. Reasonable domain.} For the pen in \#2, state the feasible domain of $A(x)$ and explain, in one sentence,
why a width outside it makes no sense. \writeline
\end{extensionbox}

\begin{spiralbox}
\textbf{That's a wrap on Unit 2!} You can now graph a parabola from any of its three forms, factor and solve by every
method (square roots, completing the square, the quadratic formula), work with complex numbers, solve systems, and
--- today --- \emph{build} a quadratic model and read the answer off its vertex and zeros. \textbf{Coming next: the
Unit 2 Review \& Test}, which pulls all of it together. Revisit your warm-ups and exit tickets as you study.
\end{spiralbox}

\end{document}
